\documentclass[12pt,a4paper,oneside]{report}

% ==============================================================================
% PACKAGE CONFIGURATIONS & PREAMBLE
% ==============================================================================
\usepackage[utf8]{inputenc}
\usepackage[T1]{fontenc}
\usepackage{geometry}
\geometry{
    top=1.0in,
    bottom=1.0in,
    left=1.0in,
    right=1.0in
}

\usepackage{amsmath,amssymb,amsfonts,amsthm}
\usepackage{graphicx}
\usepackage{booktabs}
\usepackage{tabularx}
\usepackage{multirow}
\usepackage{array}
\usepackage{listings}
\usepackage{xcolor}
\usepackage{fancyhdr}
\usepackage{titlesec}
\usepackage{setspace}
\usepackage{tikz}
\usetikzlibrary{shapes.geometric, arrows.meta, positioning, shadows, calc, fit, backgrounds}
\usepackage{algorithm}
\usepackage{algpseudocode}
\usepackage[hyphens]{url}
\usepackage[colorlinks=true, linkcolor=blue!80!black, citecolor=green!60!black, urlcolor=cyan!80!black]{hyperref}

% Typography & Spacing
\onehalfspacing
\setlength{\parindent}{0pt}
\setlength{\parskip}{6pt plus 2pt minus 1pt}

% Code Listing Styling (C++, JavaScript, SQL, Bash)
\definecolor{codegreen}{rgb}{0.13,0.55,0.13}
\definecolor{codegray}{rgb}{0.45,0.45,0.45}
\definecolor{codepurple}{rgb}{0.58,0,0.82}
\definecolor{backcolour}{rgb}{0.97,0.98,0.99}
\definecolor{codeblue}{rgb}{0.05,0.35,0.75}

\lstdefinestyle{customstyle}{
    backgroundcolor=\color{backcolour},
    commentstyle=\color{codegreen}\itshape,
    keywordstyle=\color{codeblue}\bfseries,
    numberstyle=\tiny\color{codegray},
    stringstyle=\color{codepurple},
    basicstyle=\ttfamily\footnotesize,
    breakatwhitespace=false,
    breaklines=true,
    captionpos=b,
    keepspaces=true,
    numbers=left,
    numbersep=8pt,
    showspaces=false,
    showstringspaces=false,
    showtabs=false,
    tabsize=2,
    frame=single,
    rulecolor=\color{black!15}
}
\lstset{style=customstyle}

% Header and Footer Styling
\pagestyle{fancy}
\fancyhf{}
\fancyhead[L]{\small \textbf{ParaHist:} Parallel Histogram Generation Using OpenMP}
\fancyhead[R]{\small \nouppercase{\leftmark}}
\fancyfoot[C]{\thepage}
\renewcommand{\headrulewidth}{0.4pt}
\renewcommand{\footrulewidth}{0.4pt}

% Chapter formatting
\titleformat{\chapter}[display]
  {\normalfont\huge\bfseries\color{blue!80!black}}{\chaptertitlename\ \thechapter}{14pt}{\Huge}
\titlespacing*{\chapter}{0pt}{0pt}{20pt}

% Mathematical Theorems
\newtheorem{theorem}{Theorem}[chapter]
\newtheorem{lemma}[theorem]{Lemma}
\newtheorem{definition}[theorem]{Definition}
\newtheorem{corollary}[theorem]{Corollary}

% ==============================================================================
% DOCUMENT BEGINS
% ==============================================================================
\begin{document}

% ------------------------------------------------------------------------------
% TITLE PAGE
% ------------------------------------------------------------------------------
\begin{titlepage}
    \centering
    \vspace*{0.5cm}
    
    {\Huge \textbf{ParaHist: High-Performance Parallel Histogram Generation Using OpenMP}}\\[0.6cm]
    {\LARGE \textbf{An End-to-End Scalability Analytics Engine and Distributed System Architecture}}\\[1.2cm]
    
    \rule{\linewidth}{1.5pt}\\[0.6cm]
    {\large \textbf{Academic Project Report and Comprehensive System Specification}}\\[0.3cm]
    \rule{\linewidth}{1.5pt}\\[2.0cm]
    
    \begin{minipage}{0.48\textwidth}
        \begin{flushleft} \large
            \textbf{Author:}\\
            Ajay Kumar K R\\
            \textit{Lead Systems \& HPC Researcher}\\
            Department of Computer Science
        \end{flushleft}
    \end{minipage}
    \begin{minipage}{0.48\textwidth}
        \begin{flushright} \large
            \textbf{Project Domain:}\\
            High-Performance Computing (HPC)\\
            Shared-Memory Parallelism\\
            Systems \& Web Architecture
        \end{flushright}
    \end{minipage}\\[2.5cm]
    
    \textbf{Target Architecture:}\\
    Multi-Core x86\_64 SMP $\bullet$ OpenMP 4.5/5.0 $\bullet$ C++17 $\bullet$ WSL2 Linux Kernel\\
    Node.js v20+ $\bullet$ PostgreSQL 18.6 Enterprise Database $\bullet$ React 18 SPA\\[1.5cm]
    
    \textbf{Dataset:}\\
    MNIST Handwritten Digit Corpus\\
    (42,000 Grayscale Images $\times$ 784 Pixels = 32,928,000 Pixels, 256 Intensity Bins)\\[1.5cm]
    
    \vfill
    {\large \textbf{Academic Year 2026 -- Published October 2026}}
\end{titlepage}

% ------------------------------------------------------------------------------
% PRELIMINARY PAGES: ABSTRACT & TABLE OF CONTENTS
% ------------------------------------------------------------------------------
\pagenumbering{roman}

\chapter*{Abstract}
Digital image processing, computer vision, and scientific data analytics ubiquitously depend on statistical intensity histograms for contrast adjustment, segmentation, thresholding, and probability density modeling. In data-intensive workloads involving medical scans, satellite imagery, and neural network corpora, histograms must be generated over tens or hundreds of millions of pixel evaluations.

Single-threaded computation on modern multi-core processors constitutes a severe computational bottleneck, leaving more than $85\%$ of silicon processing throughput idle. However, parallelizing histogram computation is notoriously difficult due to \textbf{write-conflict data hazards} (where concurrent threads attempt to increment identical intensity bins) and \textbf{cache-line false sharing} (where cache coherence invalidations bounce shared cache lines across cores).

\textbf{ParaHist} is an enterprise-grade, high-performance computing research application and interactive web platform that designs, implements, benchmarks, and analyzes parallel histogram generation on a shared-memory architecture. Built with \textbf{C++17 and OpenMP}, ParaHist adopts an optimal \textbf{thread-private accumulation strategy followed by a deterministic reduction}, completely eliminating locking contention and false sharing.

To make high-performance computing metrics accessible to researchers, students, and evaluators, the platform integrates:
\begin{enumerate}
    \item An optimized native C++17 OpenMP histogram and benchmark engine compiled with GCC \texttt{-O3} under the WSL2 Linux environment.
    \item An asynchronous \textbf{Node.js/Express REST API} utilizing cross-subsystem execution bridges.
    \item An enterprise-grade \textbf{PostgreSQL 18.6} relational database managing users, preferences, execution runs, 256-bin histogram distributions, and scalability benchmarks.
    \item A reactive \textbf{React 18} frontend featuring dynamic SVG/HTML5 visual analytics (Recharts), an automated product walkthrough (Driver.js), a comprehensive Account Center, and a dual-theme design system.
    \item An automated \textbf{server-side vector PDF report generator} (\texttt{PDFKit}) that synthesizes multi-page academic research reports with charts, performance tables, and execution telemetry in real time.
\end{enumerate}

Empirical validation on the standard MNIST image dataset (42,000 images $\times$ 784 pixels = \textbf{32,928,000 total pixel evaluations}, 256 intensity bins) demonstrates that ParaHist achieves:
\begin{itemize}
    \item A reduction in execution time from \textbf{43.34 ms} (single thread) down to \textbf{12.52 ms} on 4 threads.
    \item A \textbf{3.46$\times$ real-world speedup} on 4 physical threads, yielding an outstanding \textbf{86.51\% parallel efficiency}.
    \item Peak processing throughput exceeding \textbf{2.63 Billion pixels per second} (2,630 Mpixels/sec).
    \item \textbf{100.00\% numerical correctness} across all 256 bins ($\Delta = 0$ errors), verified against the sequential baseline on every run.
\end{itemize}

\clearpage
\tableofcontents
\clearpage
\listoffigures
\clearpage
\listoftables
\clearpage

% ------------------------------------------------------------------------------
% CHAPTER 1: INTRODUCTION & PROBLEM STATEMENT
% ------------------------------------------------------------------------------
\pagenumbering{arabic}
\chapter{Introduction \& Problem Statement}

\section{Context and Importance of Image Histograms}
A histogram of a digital grayscale image is a discrete function $H(b)$ that maps each intensity level $b \in [0, B-1]$ (where $B=256$ for standard 8-bit unsigned integer pixels) to the total number of pixels in the image possessing that intensity:
\begin{equation}
    H(b) = \sum_{i=0}^{N-1} \mathbb{I}(P[i] == b)
\end{equation}
where $P[i]$ denotes the value of the $i$-th pixel, $N$ is the total pixel count, and $\mathbb{I}(\cdot)$ is the indicator function returning 1 if the condition holds and 0 otherwise.

Histograms serve as the foundational backbone for:
\begin{itemize}
    \item \textbf{Contrast Enhancement:} Histogram Equalization (HE) and Contrast Limited Adaptive Histogram Equalization (CLAHE).
    \item \textbf{Segmentation:} Otsu's thresholding and multi-level image partitioning.
    \item \textbf{Feature Extraction:} Local Binary Patterns (LBP) and Histograms of Oriented Gradients (HOG) in computer vision.
    \item \textbf{Machine Learning Preprocessing:} Normalization, pixel density verification, and anomaly detection.
\end{itemize}

\section{The Computational Challenge}
When processing high-resolution image sets, medical imagery, or extensive datasets such as MNIST (32.93 million pixels), processing each pixel sequentially requires substantial CPU cycle counts. A naive loop executing 32.93 million read, index, and increment instructions introduces significant execution latency:
\begin{verbatim}
Input Pixels:  [ 0, 255, 128, 0, 42, 255, 0, ... ]  (32,928,000 elements)
Sequential:    Pixel -> Load -> Increment Count[Pixel] -> Store
Total Cycles:  ~3.3 x 10^7 iterations
\end{verbatim}

\section{The Pitfalls of Naive Parallelization}
Parallelizing this loop across multiple threads (e.g., using OpenMP or POSIX threads) introduces fundamental parallel computing hazards:
\begin{enumerate}
    \item \textbf{Data Race / Write Conflict:} If Thread 0 and Thread 1 encounter pixel value \texttt{0} at the same time, both execute \texttt{counts[0]++}. Without synchronization, an increment is lost due to interleaved read-modify-write memory operations.
    \item \textbf{Lock Overhead (Atomic / Critical):} Applying \texttt{\#pragma omp atomic} or \texttt{\#pragma omp critical} to serialize increments eliminates races but causes extreme thread contention. Threads stall waiting for cache line ownership, resulting in parallel code that runs \textbf{5$\times$ to 20$\times$ slower} than single-threaded code.
    \item \textbf{False Sharing:} If threads update adjacent memory locations that reside on the same 64-byte L1 cache line, the CPU cache coherence protocol (MESI) constantly invalidates the cache line across cores, causing catastrophic \textbf{cache line bouncing}.
\end{enumerate}

\section{Research Objectives of ParaHist}
ParaHist was designed to systematically solve these challenges. The core research objectives are:
\begin{enumerate}
    \item \textbf{Zero-Contention Architecture:} Implement a thread-private histogram allocation model where each thread accumulates counts into its own isolated memory space without locks.
    \item \textbf{Deterministic Reduction:} Execute an efficient parallel or serial reduction across thread-private arrays to produce the final 256-bin global histogram with zero synchronization overhead.
    \item \textbf{Mathematical Correctness:} Guarantee $100\%$ numerical equivalence with sequential computation across all 256 bins for all thread configurations $T \in [1, 16]$.
    \item \textbf{Comprehensive Scalability Profiling:} Measure speedup, parallel efficiency, throughput, and standard deviation over multi-run trials to isolate CPU throttling, turbo boost, and memory bus saturation effects.
    \item \textbf{Full-Stack Accessibility:} Encapsulate the HPC engine in an enterprise architecture featuring Node.js, PostgreSQL 18.6, React 18, and vector PDF reporting.
\end{enumerate}

% ------------------------------------------------------------------------------
% CHAPTER 2: THEORETICAL FOUNDATIONS OF PARALLEL COMPUTING
% ------------------------------------------------------------------------------
\chapter{Theoretical Foundations of Parallel Computing}

Understanding the performance characteristics of ParaHist requires analyzing the underlying principles of computer architecture, memory models, and parallel scalability laws.

\section{Concurrency vs. Parallelism}
\begin{itemize}
    \item \textbf{Concurrency} is the architectural property of structuring a system into multiple independent processes or threads of control that can execute in overlapping time periods (often via time-slicing on a single CPU core).
    \item \textbf{Parallelism} is the simultaneous physical execution of multiple computational tasks at the exact same instant on distinct physical processing cores. ParaHist achieves physical hardware parallelism using OpenMP thread binding across multi-core CPUs.
\end{itemize}

\section{Flynn's Taxonomy and Shared-Memory MIMD Systems}
Under Michael J. Flynn's classical taxonomy of computer architectures:
\begin{itemize}
    \item \textbf{SISD (Single Instruction, Single Data):} Traditional single-threaded Von Neumann CPU execution.
    \item \textbf{SIMD (Single Instruction, Multiple Data):} Vector extensions (e.g., Intel AVX2, AVX-512, ARM NEON).
    \item \textbf{MISD (Multiple Instruction, Single Data):} Specialized fault-tolerant or systolic arrays.
    \item \textbf{MIMD (Multiple Instruction, Multiple Data):} Modern multi-core symmetric multiprocessing (SMP) and non-uniform memory access (NUMA) systems.
\end{itemize}

ParaHist targets \textbf{Shared-Memory MIMD} architectures. Each processing core executes an independent stream of instructions (fetching distinct pixel ranges), operating on distinct memory locations within a globally shared address space.

\section{Amdahl's Law and Gustafson's Law}

\subsection{Amdahl's Law (Strong Scaling)}
Amdahl's Law models the theoretical speedup achievable by a parallel program under a fixed problem size when executed across $p$ processing units. If $f$ is the fraction of the computation that is strictly sequential (cannot be parallelized), and $(1 - f)$ is the parallelizable fraction, the maximum speedup $S(p)$ is bounded by:
\begin{equation}
    S(p) = \frac{T_1}{T_p} = \frac{1}{f + \frac{1 - f}{p}}
\end{equation}
As the number of processors $p \to \infty$, the maximum theoretical speedup asymptotically approaches:
\begin{equation}
    \lim_{p \to \infty} S(p) = \frac{1}{f}
\end{equation}

In ParaHist:
\begin{itemize}
    \item The parallelizable component $(1 - f)$ is the iteration through the $32,928,000$ pixels and local accumulation.
    \item The sequential component $f$ includes thread initialization, OpenMP fork/join barrier overhead, and the final reduction loop merging $p \times 256$ bins into the global array.
    \item For 4 threads with $f \approx 0.04$ (4\% serial overhead):
    \begin{equation}
        S(4) = \frac{1}{0.04 + \frac{0.96}{4}} = \frac{1}{0.04 + 0.24} = \frac{1}{0.28} \approx 3.57\times
    \end{equation}
    This theoretical bound closely matches ParaHist's empirically observed \textbf{$3.46\times$ speedup}.
\end{itemize}

\subsection{Gustafson's Law (Weak Scaling)}
When problem size scales proportionally with the number of processors, Gustafson's Law models scaled speedup:
\begin{equation}
    S_v(p) = p - f(p - 1)
\end{equation}
This highlights that as dataset sizes expand into gigabytes, the relative impact of the sequential merge diminishes further.

\section{Hardware Hierarchy, Cache Architecture, and the Memory Wall}
Modern x86\_64 microprocessors feature a tiered memory hierarchy:
\begin{enumerate}
    \item \textbf{L1 Data Cache:} 32 KB to 64 KB per core, latency $\sim 1$ ns (4--5 clock cycles).
    \item \textbf{L2 Unified Cache:} 512 KB to 1 MB per core, latency $\sim 3$--$4$ ns (12--14 clock cycles).
    \item \textbf{L3 Shared Cache:} 16 MB to 64 MB shared across cores, latency $\sim 10$--$15$ ns (35--40 clock cycles).
    \item \textbf{Main Memory (DRAM):}\ 16 GB to 64 GB DDR4/DDR5, latency $\sim 60$--$80$ ns (150--200 clock cycles).
\end{enumerate}

A critical architectural consideration in histogram generation is \textbf{bandwidth saturation (The Memory Wall)}. In ParaHist, each pixel is an 8-bit unsigned integer (\texttt{uint8\_t}), totaling $32.93\text{ MB}$ of input data. While reading $32.93\text{ MB}$ sequentially is bandwidth-efficient due to CPU hardware prefetchers, writing to memory requires cache-line linefills. By allocating each thread's local histogram to fit entirely within that core's L1 cache ($256 \times 8\text{ bytes} = 2,048\text{ bytes} = 2\text{ KB} \ll 32\text{ KB}$), \textbf{all thread histogram increments execute at zero-wait L1 cache latency}.

\section{The False Sharing Phenomenon and Cache Invalidation (MESI Protocol)}
Multi-core processors maintain memory consistency across CPU caches using cache coherence protocols such as \textbf{MESI} (Modified, Exclusive, Shared, Invalid):
\begin{itemize}
    \item Cache lines are typically 64 bytes wide (capable of holding eight 64-bit integers).
    \item If Thread 0 writes to \texttt{counts[0]} and Thread 1 writes to \texttt{counts[1]}, both variables share the identical 64-byte physical cache line.
    \item When Core 0 writes to \texttt{counts[0]}, Core 0's cache controller marks the cache line as \textbf{Modified} and transmits an invalidation signal via the processor interconnect (QPI/UPI or Infinity Fabric).
    \item Core 1's cache line is forced to the \textbf{Invalid} state. When Core 1 attempts to write to \texttt{counts[1]}, a cache miss occurs. Core 1 must stall while reloading the cache line from L3 cache or RAM.
    \item This continuous invalidation loop---termed \textbf{False Sharing} or \textbf{Cache Ping-Pong}---degrades multi-threaded performance by an order of magnitude.
\end{itemize}

\textbf{ParaHist's Architectural Mitigation:}  
ParaHist allocates independent, isolated arrays for each thread. Each thread's histogram array occupies $2,048$ bytes ($32$ distinct cache lines), far exceeding the 64-byte threshold. As a result, no two threads ever access or invalidate the same cache line during parallel execution.

\section{Synchronization Paradigms Comparison}
Table~\ref{tab:sync_comparison} summarizes the three primary synchronization approaches:

\begin{table}[h!]
\centering
\caption{Comparison of Parallel Histogram Synchronization Approaches}
\label{tab:sync_comparison}
\begin{tabularx}{\textwidth}{lXXXl}
\toprule
\textbf{Strategy} & \textbf{Synchronization Overhead} & \textbf{Race Hazard} & \textbf{False Sharing Risk} & \textbf{Scalability Factor} \\
\midrule
Global Array + Mutex & Catastrophic & None & High & $<0.05\times$ (Degrades) \\
Global Array + Atomic & Severe (\texttt{LOCK XADD}) & None & Severe & $\sim 0.3\times$--$0.8\times$ (Sub-linear) \\
\textbf{Thread-Private Buffers} & \textbf{Zero during compute} & \textbf{None} & \textbf{Zero} & \textbf{$>3.4\times$ on 4 Cores} \\
\bottomrule
\end{tabularx}
\end{table}

% ------------------------------------------------------------------------------
% CHAPTER 3: SYSTEM ARCHITECTURE & MULTI-TIER TOPOLOGY
% ------------------------------------------------------------------------------
\chapter{System Architecture \& Multi-Tier Topology}

ParaHist implements a decoupled 4-tier architecture designed for high throughput, memory safety, and cross-platform flexibility.

\section{Global 4-Tier Architectural Topology}

\begin{figure}[h!]
\centering
\begin{tikzpicture}[
    node distance=1.8cm and 1.5cm,
    tierbox/.style={draw=blue!70!black, fill=blue!5, very thick, rounded corners=8pt, inner sep=12pt},
    compbox/.style={draw=black!70, fill=white, thick, rounded corners=4pt, align=center, font=\small, minimum height=1.0cm, minimum width=2.6cm},
    arrow/.style={-{Stealth[length=3mm]}, thick, draw=blue!80!black}
]

% Tier 1: Client
\node[compbox] (ui) {React 18 SPA\\Interactive UI};
\node[compbox, right=0.6cm of ui] (chart) {Recharts 256\\Bin Visualizer};
\node[compbox, right=0.6cm of chart] (tour) {Guided Tour\\Account Center};
\begin{scope}[on background layer]
    \node[tierbox, fit=(ui)(chart)(tour), label=above:{\textbf{\large Tier 1: Presentation Layer (React 18 / Vite / Tailwind)}}] (tier1) {};
\end{scope}

% Tier 2: Backend
\node[compbox, below=1.8cm of chart] (express) {Node.js / Express\\API Gateway (:5000)};
\node[compbox, left=0.6cm of express] (auth) {Auth Middleware\\\& JWT Engine};
\node[compbox, right=0.6cm of express] (pdf) {PDFKit Vector\\Report Builder};
\begin{scope}[on background layer]
    \node[tierbox, fit=(auth)(express)(pdf), fill=green!5, draw=green!70!black, label=above:{\textbf{\large Tier 2: Application API Gateway (Node.js Express)}}] (tier2) {};
\end{scope}

% Tier 3: HPC Layer
\node[compbox, below=1.8cm of auth] (wsl) {WSL2 Ubuntu\\Subsystem Bridge};
\node[compbox, below=1.8cm of express] (cpp) {C++17 OpenMP Engine\\(\texttt{parahist\_app})};
\node[compbox, below=1.8cm of pdf] (mnist) {MNIST Dataset\\(32.93M Pixels)};
\begin{scope}[on background layer]
    \node[tierbox, fit=(wsl)(cpp)(mnist), fill=orange!5, draw=orange!70!black, label=above:{\textbf{\large Tier 3: HPC Computation Layer (WSL2 / C++17 OpenMP)}}] (tier3) {};
\end{scope}

% Tier 4: Database Layer
\node[compbox, below=1.8cm of cpp] (pg) {PostgreSQL 18.6 Relational Database\\(users, experiments, 256-bin histograms, benchmark\_runs)};
\begin{scope}[on background layer]
    \node[tierbox, fit=(pg), fill=purple!5, draw=purple!70!black, label=above:{\textbf{\large Tier 4: Relational Persistence Layer (PostgreSQL 18.6)}}] (tier4) {};
\end{scope}

% Inter-tier Connections
\draw[arrow] (chart) -- node[right, font=\scriptsize] {REST / JSON} (express);
\draw[arrow] (express) -- node[left, font=\scriptsize] {child\_process} (wsl);
\draw[arrow] (wsl) -- (cpp);
\draw[arrow] (mnist) -- (cpp);
\draw[arrow] (express) -- node[right, font=\scriptsize] {SQL Queries / pg pool} (pg);

\end{tikzpicture}
\caption{Global 4-Tier System Architecture of the ParaHist Platform}
\label{fig:system_arch}
\end{figure}

\section{Cross-Platform Execution Bridge (Windows 11 to WSL2)}
High-performance C++ OpenMP code relies heavily on \texttt{gcc}, \texttt{libgomp}, and native POSIX thread scheduling for deterministic timing. Windows OS thread scheduling introduces background preemption and non-standard thread affinities. To solve this without requiring the user to switch operating systems, ParaHist features a \textbf{Cross-Subsystem Virtualization Bridge}:
\begin{enumerate}
    \item \textbf{Host Environment:} Windows 11 running Node.js v20.x and PostgreSQL 18.6 on \texttt{localhost:5432}.
    \item \textbf{Virtual Execution Environment:} Windows Subsystem for Linux (WSL2) hosting Ubuntu 22.04 LTS.
    \item \textbf{Execution Invocation:} The Node.js application executes native Linux binaries inside WSL using \texttt{child\_process.execFile}:
    \begin{verbatim}
wsl.exe -d Ubuntu-22.04 -- /home/user/parahist/build/parahist_app --parallel 8
    \end{verbatim}
    \item \textbf{Inter-Process Communication:} The native binary streams structured JSON over standard output (\texttt{stdout}), which is captured by the Node.js stream buffer, validated, parsed, and committed to PostgreSQL in an atomic transaction.
\end{enumerate}

% ------------------------------------------------------------------------------
% CHAPTER 4: ALGORITHMIC DESIGN & MATHEMATICAL FORMULATIONS
% ------------------------------------------------------------------------------
\chapter{Algorithmic Design \& Mathematical Formulations}

\section{Problem Mathematical Formulation}
Given an image dataset $\mathcal{D} = \{R_0, R_1, \dots, R_{M-1}\}$ consisting of $M = 42,000$ image rows, where each row $R_i$ contains $K = 784$ discrete pixel samples:
\begin{equation}
    R_i = [p_{i,0}, p_{i,1}, \dots, p_{i,K-1}], \quad p_{i,j} \in \{0, 1, \dots, 255\}
\end{equation}
The total input domain consists of:
\begin{equation}
    N = M \times K = 42,000 \times 784 = 32,928,000 \text{ pixels}
\end{equation}
The objective is to compute the global histogram vector $\mathbf{H} \in \mathbb{N}^{256}$:
\begin{equation}
    \mathbf{H}[b] = \sum_{i=0}^{M-1} \sum_{j=0}^{K-1} [p_{i,j} == b], \quad \forall b \in \{0, 1, \dots, 255\}
\end{equation}
satisfying the conservation property:
\begin{equation}
    \sum_{b=0}^{255} \mathbf{H}[b] = N = 32,928,000
\end{equation}

\section{Algorithm 1: Baseline Sequential Histogram Computation}

\begin{algorithm}[h!]
\caption{Sequential Pixel Histogram Generation}
\label{alg:sequential}
\begin{algorithmic}[1]
\Require Dataset rows $R$ of size $M$, where each row has $K$ pixels in $[0, 255]$
\Ensure HistogramResult containing $counts$ array of size 256 and $elapsed\_ms$
\State $counts \gets \text{Array of 256 elements, all initialized to } 0$
\State $t_{start} \gets \text{HighResolutionTimer::Now}()$
\For{each row $r \in R$}
    \For{each pixel $px \in r.pixels$}
        \State $counts[px] \gets counts[px] + 1$
    \EndFor
\EndFor
\State $t_{end} \gets \text{HighResolutionTimer::Now}()$
\State $elapsed\_ms \gets (t_{end} - t_{start}) \text{ in milliseconds}$
\State \Return $\{ counts, elapsed\_ms \}$
\end{algorithmic}
\end{algorithm}

\section{Algorithm 2: Parallel Histogram via Thread-Private Reduction}

\begin{algorithm}[h!]
\caption{Parallel Histogram via Thread-Private Reduction}
\label{alg:parallel}
\begin{algorithmic}[1]
\Require Dataset rows $R$ of size $M$, Thread count $P \in [1, 16]$
\Ensure HistogramResult containing $counts$ array of size 256 and $elapsed\_ms$
\State $nthreads \gets (P > 0) \text{ ? } P : \text{omp\_get\_max\_threads}()$
\State $local\_hists \gets \text{Matrix of size } [nthreads][256] \text{ zero-initialized}$
\State $t_{start} \gets \text{omp\_get\_wtime}()$
\Statex
\State \textbf{\#pragma omp parallel num\_threads(nthreads)}
\State \Begin
    \State $tid \gets \text{omp\_get\_thread\_num}()$
    \State $local \gets \text{Reference to } local\_hists[tid]$
    \Statex \Comment{Static scheduling assigns contiguous chunks of size $M/nthreads$}
    \State \textbf{\#pragma omp for schedule(static)}
    \For{$i \gets 0 \text{ to } M - 1$}
        \For{each pixel $px \in R[i].pixels$}
            \State $local[px] \gets local[px] + 1$ \Comment{Write exclusively to private array}
        \EndFor
    \EndFor
\State \End \Comment{Implicit barrier: all $P$ threads join here}
\Statex
\State $final\_counts \gets \text{Array of 256 elements initialized to } 0$
\For{$t \gets 0 \text{ to } nthreads - 1$}
    \For{$b \gets 0 \text{ to } 255$}
        \State $final\_counts[b] \gets final\_counts[b] + local\_hists[t][b]$
    \EndFor
\EndFor
\State $t_{end} \gets \text{omp\_get\_wtime}()$
\State $elapsed\_ms \gets (t_{end} - t_{start}) \times 1000.0$
\State \Return $\{ final\_counts, elapsed\_ms \}$
\end{algorithmic}
\end{algorithm}

\section{Theoretical Complexity \& Correctness Proof}

\subsection{Complexity Analysis}
\begin{itemize}
    \item \textbf{Parallel Counting Phase:} Each thread evaluates $\frac{M \times K}{P} = \frac{N}{P}$ pixel values.
    \item \textbf{Reduction Phase:} Merging $P$ thread arrays requires $P \times 256$ operations.
    \item \textbf{Total Time Complexity:}
    \begin{equation}
        T(P) = \mathcal{O}\left(\frac{N}{P} + P \times B\right)
    \end{equation}
    Since $N = 32,928,000$ and $B = 256$, for $P=4$: $\frac{32,928,000}{4} = 8,232,000 \gg 1,024$. The reduction represents $<0.012\%$ of total execution time.
\end{itemize}

\subsection{Proof of Race-Freedom and Equivalence}
\begin{lemma}[Absence of Data Races]
During the parallel region of Algorithm~\ref{alg:parallel}, no memory location is written to by more than one thread.
\end{lemma}
\begin{proof}
Each thread writes strictly to $local\_hists[tid]$. Since thread IDs $tid \in [0, P-1]$ are unique and their buffer memory regions do not overlap, Bernstein's condition $O_u \cap O_v = \emptyset$ holds for all $u \neq v$. Thus, the parallel section is provably free of write conflicts.
\end{proof}

\begin{theorem}[Numerical Equivalence]
The final parallel histogram $\mathbf{H}_{par}$ equals the sequential histogram $\mathbf{H}_{seq}$ for all $b \in [0, 255]$.
\end{theorem}
\begin{proof}
The static schedule partitions image indices $\{0, \dots, M-1\}$ into disjoint subsets $\bigcup_{t=0}^{P-1} \mathcal{I}_t$. By associativity of integer addition:
\begin{equation}
    \mathbf{H}_{par}[b] = \sum_{t=0}^{P-1} \sum_{i \in \mathcal{I}_t} \sum_{j=0}^{K-1} [p_{i,j} == b] = \sum_{i=0}^{M-1} \sum_{j=0}^{K-1} [p_{i,j} == b] = \mathbf{H}_{seq}[b]
\end{equation}
Hence, numerical equivalence is preserved identically.
\end{proof}

% ------------------------------------------------------------------------------
% CHAPTER 5: DETAILED FLOW DIAGRAMS & EXECUTION WORKFLOWS
% ------------------------------------------------------------------------------
\chapter{Detailed Flow Diagrams \& Execution Workflows}

\section{C++ OpenMP Parallel Execution Engine Flowchart}

\begin{figure}[h!]
\centering
\begin{tikzpicture}[
    node distance=1.2cm,
    flowstep/.style={draw=blue!70!black, fill=blue!5, thick, rounded corners=3pt, minimum width=4.5cm, minimum height=0.8cm, align=center, font=\small},
    decision/.style={diamond, draw=red!70!black, fill=red!5, thick, aspect=2, align=center, font=\small},
    arrow/.style={-{Stealth[length=2.5mm]}, thick, draw=blue!80!black}
]

\node[flowstep] (start) {Process Spawned: \texttt{parahist\_app --parallel N}};
\node[flowstep, below of=start] (load) {Ingest MNIST CSV ($42,000 \times 784$ Pixels)};
\node[flowstep, below of=load] (alloc) {Allocate $local\_hists[N][256]$ Zero-Initialized};
\node[flowstep, below of=alloc] (timer1) {Record Wall Clock Start: \texttt{omp\_get\_wtime()}};
\node[flowstep, below of=timer1] (fork) {\texttt{\#pragma omp parallel num\_threads(N)}};
\node[flowstep, below of=fork] (static) {Partition Rows Contiguously: \texttt{schedule(static)}};
\node[flowstep, below of=static] (count) {Thread-Private Counting: \texttt{local[pixel]++}};
\node[flowstep, below of=count] (barrier) {Implicit OpenMP Join Barrier};
\node[flowstep, below of=barrier] (merge) {Serial Merge into Global Histogram};
\node[flowstep, below of=merge] (timer2) {Record Wall Clock Stop: \texttt{omp\_get\_wtime()}};
\node[flowstep, below of=timer2] (json) {Format Output JSON \& Exit Code 0};

\draw[arrow] (start) -- (load);
\draw[arrow] (load) -- (alloc);
\draw[arrow] (alloc) -- (timer1);
\draw[arrow] (timer1) -- (fork);
\draw[arrow] (fork) -- (static);
\draw[arrow] (static) -- (count);
\draw[arrow] (count) -- (barrier);
\draw[arrow] (barrier) -- (merge);
\draw[arrow] (merge) -- (timer2);
\draw[arrow] (timer2) -- (json);

\end{tikzpicture}
\caption{Flowchart of C++17 OpenMP Thread-Private Histogram Engine}
\label{fig:cpp_flowchart}
\end{figure}

\section{Full-Stack Request/Response Execution Flowchart}
Figure~\ref{fig:fullstack_flowchart} illustrates the complete lifecycle from the browser user click to database persistence:

\begin{figure}[h!]
\centering
\begin{tikzpicture}[
    node distance=1.1cm,
    block/.style={draw=black!70, fill=gray!5, thick, rounded corners=3pt, minimum width=6cm, minimum height=0.7cm, align=center, font=\small},
    arrow/.style={-{Stealth[length=2.5mm]}, thick, draw=blue!80!black}
]

\node[block] (s1) {User selects thread count \& clicks \textbf{"Run Experiment"}};
\node[block, below of=s1] (s2) {Frontend dispatches \texttt{POST /api/histogram} with JWT};
\node[block, below of=s2] (s3) {Express Router verifies Bearer token \& sanitizes inputs};
\node[block, below of=s3] (s4) {\texttt{cppRunner.js} executes C++ binary inside WSL2 Ubuntu};
\node[block, below of=s4] (s5) {WSL reads \texttt{train.csv} \& runs OpenMP calculation};
\node[block, below of=s5] (s6) {JSON response parsed from process \texttt{stdout}};
\node[block, below of=s6] (s7) {PostgreSQL commits experiment metrics \& 256 bins};
\node[block, below of=s7] (s8) {Express delivers HTTP 200 payload to React frontend};
\node[block, below of=s8] (s9) {Recharts dynamically animates 256-bin histogram curve};

\draw[arrow] (s1) -- (s2);
\draw[arrow] (s2) -- (s3);
\draw[arrow] (s3) -- (s4);
\draw[arrow] (s4) -- (s5);
\draw[arrow] (s5) -- (s6);
\draw[arrow] (s6) -- (s7);
\draw[arrow] (s7) -- (s8);
\draw[arrow] (s8) -- (s9);

\end{tikzpicture}
\caption{End-to-End Application Data Flow Lifecycle}
\label{fig:fullstack_flowchart}
\end{figure}

% ------------------------------------------------------------------------------
% CHAPTER 6: DATABASE ARCHITECTURE & RELATIONAL SCHEMA (POSTGRESQL 18.6)
% ------------------------------------------------------------------------------
\chapter{Database Architecture \& Relational Schema (PostgreSQL 18.6)}

ParaHist integrates \textbf{PostgreSQL 18.6} as its primary persistent relational data store. This replaces legacy file-based JSON storage, ensuring ACID transaction compliance, strict foreign key constraints, high-concurrency connection pooling, and optimized query execution plans.

\section{Entity-Relationship (ER) Architecture}
The database schema consists of six primary relational entities:
\begin{enumerate}
    \item \texttt{users}: Master account entity storing authenticated user records.
    \item \texttt{user\_preferences}: One-to-one preferences (theme, default thread count, histogram mode).
    \item \texttt{notification\_preferences}: One-to-one alert toggles.
    \item \texttt{experiments}: Historical trial records capturing sequential/parallel times and speedup.
    \item \texttt{experiment\_histogram}: Normalized 256-row intensity distributions per experiment.
    \item \texttt{benchmark\_runs}: Multi-run scalability benchmarks and variance statistics.
\end{enumerate}

\section{PostgreSQL DDL Schema Implementation}
The following SQL defines the database structure idempotently:

\begin{lstlisting}[language=SQL, caption={PostgreSQL 18.6 Schema Definition for ParaHist}]
-- 1. Table: users
CREATE TABLE IF NOT EXISTS users (
  id VARCHAR(64) PRIMARY KEY,
  name VARCHAR(255) NOT NULL,
  email VARCHAR(255) UNIQUE NOT NULL,
  password_hash VARCHAR(255) NOT NULL,
  role VARCHAR(50) DEFAULT 'Student',
  created_at TIMESTAMPTZ DEFAULT NOW(),
  updated_at TIMESTAMPTZ DEFAULT NOW()
);

-- 2. Table: user_preferences
CREATE TABLE IF NOT EXISTS user_preferences (
  id SERIAL PRIMARY KEY,
  user_id VARCHAR(64) UNIQUE NOT NULL REFERENCES users(id) ON DELETE CASCADE,
  theme VARCHAR(20) DEFAULT 'dark',
  default_threads INTEGER DEFAULT 8,
  default_histogram_mode VARCHAR(20) DEFAULT 'both',
  guided_tour BOOLEAN DEFAULT true,
  created_at TIMESTAMPTZ DEFAULT NOW(),
  updated_at TIMESTAMPTZ DEFAULT NOW()
);

-- 3. Table: experiments
CREATE TABLE IF NOT EXISTS experiments (
  id SERIAL PRIMARY KEY,
  user_id VARCHAR(64) REFERENCES users(id) ON DELETE SET NULL,
  threads INTEGER NOT NULL,
  sequential_ms DOUBLE PRECISION,
  parallel_ms DOUBLE PRECISION,
  speedup DOUBLE PRECISION,
  efficiency DOUBLE PRECISION,
  correctness BOOLEAN DEFAULT true,
  mismatched_bins INTEGER DEFAULT 0,
  sequential_total BIGINT,
  parallel_total BIGINT,
  expected_total BIGINT DEFAULT 32928000,
  bin_count INTEGER DEFAULT 256,
  most_frequent_pixel INTEGER,
  most_frequent_count BIGINT,
  created_at TIMESTAMPTZ DEFAULT NOW()
);

-- 4. Table: experiment_histogram (256 bins per experiment)
CREATE TABLE IF NOT EXISTS experiment_histogram (
  id BIGSERIAL PRIMARY KEY,
  experiment_id INTEGER NOT NULL REFERENCES experiments(id) ON DELETE CASCADE,
  pixel_value INTEGER NOT NULL CHECK (pixel_value >= 0 AND pixel_value <= 255),
  sequential_count BIGINT NOT NULL,
  parallel_count BIGINT NOT NULL
);

-- 5. Indexes for fast query planning
CREATE INDEX IF NOT EXISTS idx_users_email ON users(email);
CREATE INDEX IF NOT EXISTS idx_experiments_user_id ON experiments(user_id);
CREATE INDEX IF NOT EXISTS idx_experiments_created_at ON experiments(created_at DESC);
CREATE INDEX IF NOT EXISTS idx_experiment_histogram_exp_id ON experiment_histogram(experiment_id);
\end{lstlisting}

% ------------------------------------------------------------------------------
% CHAPTER 7: SOURCE CODE IMPLEMENTATION & WALKTHROUGH
% ------------------------------------------------------------------------------
\chapter{Source Code Implementation \& Walkthrough}

\section{Core C++17 OpenMP Engine}

\begin{lstlisting}[language=C++, caption={C++17 OpenMP Parallel Kernel (\texttt{parallel\_histogram.cpp})}]
#include "parallel_histogram.h"
#include <omp.h>
#include <vector>
#include <array>

HistogramResult ParallelHistogram::compute(const std::vector<MnistRow>& rows,
                                            int num_threads)
{
    int nthreads = (num_threads > 0) ? num_threads : omp_get_max_threads();

    // Allocate 2D buffer: nthreads isolated 256-bin arrays
    using LocalHist = std::array<long long, HIST_BINS>;
    std::vector<LocalHist> local_hists(nthreads);
    for (auto& h : local_hists) h.fill(0LL);

    const std::size_t n_rows = rows.size();
    double t_start = omp_get_wtime();

    #pragma omp parallel num_threads(nthreads)
    {
        int tid = omp_get_thread_num();
        LocalHist& local = local_hists[tid];

        #pragma omp for schedule(static)
        for (std::size_t i = 0; i < n_rows; ++i) {
            for (const uint8_t px : rows[i].pixels) {
                local[px]++;
            }
        }
    } // Implicit barrier: All threads join before merge

    // Deterministic serial reduction on master thread
    HistogramResult result;
    result.counts.fill(0LL);

    for (int t = 0; t < nthreads; ++t) {
        for (int b = 0; b < HIST_BINS; ++b) {
            result.counts[b] += local_hists[t][b];
        }
    }

    double t_end = omp_get_wtime();
    result.elapsed_ms = (t_end - t_start) * 1000.0;
    return result;
}
\end{lstlisting}

\section{Subprocess Execution Bridge}

\begin{lstlisting}[language=JavaScript, caption={Cross-Subsystem WSL Bridge (\texttt{backend/src/services/cppRunner.js})}]
const { execFile } = require('child_process');
const path = require('path');

function runWslCommand(binaryPath, args = []) {
  return new Promise((resolve, reject) => {
    // Translate Windows paths to WSL mount paths
    const wslBinary = binaryPath.replace(/\\/g, '/')
      .replace(/^([A-Za-z]):/, (_, d) => `/mnt/${d.toLowerCase()}`);
    const wslArgs = ['-d', 'Ubuntu-22.04', '--', wslBinary, ...args];

    execFile('wsl.exe', wslArgs, { maxBuffer: 10 * 1024 * 1024 }, (error, stdout, stderr) => {
      if (error) {
        return reject(new Error(`WSL execution failed: ${stderr || error.message}`));
      }
      try {
        resolve(JSON.parse(stdout));
      } catch (e) {
        resolve({ raw: stdout });
      }
    });
  });
}
\end{lstlisting}

% ------------------------------------------------------------------------------
% CHAPTER 8: EXPERIMENTAL METHODOLOGY & BENCHMARK RESULTS
% ------------------------------------------------------------------------------
\chapter{Experimental Methodology \& Benchmark Results}

\section{Testbed Environment Specifications}
Table~\ref{tab:testbed} specifies the testbed environment used for empirical validation:

\begin{table}[h!]
\centering
\caption{System Hardware and Software Configuration}
\label{tab:testbed}
\begin{tabularx}{\textwidth}{lX}
\toprule
\textbf{Component} & \textbf{Specification} \\
\midrule
CPU Model & AMD Ryzen / Intel Core x86\_64 Multi-Core Architecture \\
Cores / Threads & 8 Physical Cores / 16 Logical SMT Threads \\
Clock Frequency & 3.20 GHz Base Clock / 4.40 GHz Max Boost Clock \\
CPU Caches & 32 KB L1d per core, 512 KB L2 per core, 16 MB L3 Unified \\
System Memory & 16 GB DDR4 Dual-Channel @ 3200 MT/s \\
Operating System & Microsoft Windows 11 Professional (Build 22631) \\
Linux Runtime & WSL2 (Kernel 5.15.153.1-microsoft-standard-WSL2) \\
C++ Compiler & GCC / G++ 11.4.0 (\texttt{-std=c++17 -O3 -fopenmp -Wall}) \\
Database & PostgreSQL 18.6 Enterprise Database \\
\bottomrule
\end{tabularx}
\end{table}

\section{Multi-Run Empirical Benchmark Data}
Table~\ref{tab:multi_run} documents the verified empirical results across 5 consecutive trials on the 32,928,000 pixel MNIST dataset:

\begin{table}[h!]
\centering
\caption{Empirical Multi-Run Benchmark Results (32,928,000 Pixels)}
\label{tab:multi_run}
\begin{tabular}{cccccc}
\toprule
\textbf{Threads ($T$)} & \textbf{Run \#} & \textbf{Sequential ($T_{seq}$)} & \textbf{Parallel ($T_{par}$)} & \textbf{Speedup ($S_T$)} & \textbf{Efficiency ($E_T$)} \\
\midrule
1 & 1 & 43.19 ms & 42.80 ms & 1.01$\times$ & 100.90\% \\
1 & 2 & 42.99 ms & 43.21 ms & 0.99$\times$ & 99.49\% \\
1 & 3 & 43.70 ms & 43.11 ms & 1.01$\times$ & 101.39\% \\
1 & 4 & 43.45 ms & 43.10 ms & 1.01$\times$ & 100.81\% \\
1 & 5 & 43.39 ms & 43.27 ms & 1.00$\times$ & 100.27\% \\
\midrule
2 & 1 & 43.09 ms & 22.46 ms & 1.92$\times$ & 95.90\% \\
2 & 2 & 43.13 ms & 22.62 ms & 1.91$\times$ & 95.33\% \\
2 & 3 & 43.88 ms & 22.33 ms & 1.97$\times$ & 98.28\% \\
2 & 4 & 43.11 ms & 22.10 ms & 1.95$\times$ & 97.55\% \\
2 & 5 & 44.78 ms & 23.45 ms & 1.91$\times$ & 95.47\% \\
\midrule
4 & 1 & 44.49 ms & 12.14 ms & 3.66$\times$ & 91.61\% \\
4 & 2 & 44.80 ms & 12.62 ms & 3.55$\times$ & 88.75\% \\
4 & 3 & 44.92 ms & 12.44 ms & 3.61$\times$ & 90.30\% \\
4 & 4 & 45.13 ms & 13.10 ms & 3.45$\times$ & 86.14\% \\
4 & 5 & 44.76 ms & 12.32 ms & 3.63$\times$ & 90.83\% \\
\bottomrule
\end{tabular}
\end{table}

\section{Scalability Summary and Throughput Metrics}
Table~\ref{tab:summary_stats} summarizes the mean performance characteristics:

\begin{table}[h!]
\centering
\caption{Aggregated Performance and Throughput Metrics}
\label{tab:summary_stats}
\begin{tabularx}{\textwidth}{ccccccc}
\toprule
\textbf{Threads} & \textbf{Mean Seq} & \textbf{Mean Par} & \textbf{Std Dev ($\sigma$)} & \textbf{Speedup} & \textbf{Efficiency} & \textbf{Throughput} \\
\midrule
\textbf{1 Thread} & 43.34 ms & 43.10 ms & 0.18 ms & \textbf{1.01$\times$} & \textbf{100.56\%} & 763.99 Mpixels/sec \\
\textbf{2 Threads} & 43.34 ms & 22.59 ms & 0.52 ms & \textbf{1.92$\times$} & \textbf{95.91\%} & 1,457.64 Mpixels/sec \\
\textbf{4 Threads} & 43.34 ms & 12.52 ms & 0.37 ms & \textbf{3.46$\times$} & \textbf{86.51\%} & \textbf{2,629.42 Mpixels/sec} \\
\bottomrule
\end{tabularx}
\end{table}

% ------------------------------------------------------------------------------
% CHAPTER 9: EVALUATOR EXECUTION GUIDE & CONCLUSION
% ------------------------------------------------------------------------------
\chapter{Evaluator Execution Guide \& Conclusion}

\section{Presentation Execution Instructions}
To run the full stack for presentation or academic evaluation:

\begin{enumerate}
    \item \textbf{Start the Backend API Gateway (Terminal 1):}
    \begin{lstlisting}[language=bash]
cd c:\Users\ajayk\OneDrive\Desktop\Parallel-mini-project\backend
npm start
    \end{lstlisting}
    \textit{Expected Console Output:}
    \begin{verbatim}
[initDatabase] Checking/verifying ParaHist PostgreSQL schema...
✓ ParaHist PostgreSQL schema and indexes verified successfully.
Server running on port 5000
Connected to PostgreSQL database: parahist on localhost:5432
    \end{verbatim}

    \item \textbf{Start the React Frontend Development Server (Terminal 2):}
    \begin{lstlisting}[language=bash]
cd c:\Users\ajayk\OneDrive\Desktop\Parallel-mini-project\frontend
npm run dev
    \end{lstlisting}
    \textit{Expected Console Output:}
    \begin{verbatim}
VITE v5.4.19 ready in 437 ms
➜  Local:   http://localhost:5173/
    \end{verbatim}

    \item \textbf{Access the Web Application:}
    Navigate to \texttt{http://localhost:5173/}. Sign in using standard evaluation credentials:
    \begin{itemize}
        \item \textbf{Email:} \texttt{demo@parahist.edu}
        \item \textbf{Password:} \texttt{demo1234}
    \end{itemize}
\end{enumerate}

\section{Conclusion}
The \textbf{ParaHist} project demonstrates the design, optimization, and empirical evaluation of high-performance parallel histogram generation. By analyzing the root causes of multi-threaded degradation---specifically \textbf{write-conflict data hazards} and \textbf{cache-line false sharing}---ParaHist implements a \textbf{thread-private accumulation strategy} that delivers lock-free, race-free parallel execution.

Key outcomes achieved include:
\begin{itemize}
    \item \textbf{3.46$\times$ Speedup on 4 Threads:} Reduces processing time from 43.34 ms down to 12.52 ms on the 32.93-million-pixel MNIST dataset.
    \item \textbf{High Parallel Efficiency:} Achieves \textbf{86.51\% efficiency}, closely matching Amdahl's theoretical upper bound.
    \item \textbf{2.63 Billion Pixels/Sec Throughput:} Delivers high processing rates with \textbf{100.00\% numerical correctness} across all 256 bins.
    \item \textbf{Production-Ready System:} Backed by an ACID-compliant PostgreSQL 18.6 database, a modern React 18 frontend, and automated vector PDF reporting.
\end{itemize}

% ------------------------------------------------------------------------------
% BIBLIOGRAPHY
% ------------------------------------------------------------------------------
\begin{thebibliography}{99}
\bibitem{openmp} OpenMP Architecture Review Board. \textit{OpenMP Application Programming Interface Specification, Version 5.0.} November 2018.
\bibitem{amdahl} Amdahl, Gene M. "Validity of the single processor approach to achieving large scale computing capabilities." \textit{AFIPS Conference Proceedings}, vol. 30, 1967, pp. 483--485.
\bibitem{gustafson} Gustafson, John L. "Reevaluating Amdahl's Law." \textit{Communications of the ACM}, vol. 31, no. 5, 1988, pp. 532--533.
\bibitem{flynn} Flynn, Michael J. "Some Computer Organizations and Their Effectiveness." \textit{IEEE Transactions on Computers}, vol. C-21, no. 9, 1972, pp. 948--960.
\bibitem{patterson} Hennessy, John L., and David A. Patterson. \textit{Computer Architecture: A Quantitative Approach.} 6th Edition, Morgan Kaufmann / Elsevier, 2017.
\bibitem{mnist} LeCun, Yann, Corinna Cortes, and Christopher J.C. Burges. "The MNIST Database of Handwritten Digits." \textit{Courant Institute \& Bell Labs}, 1998.
\bibitem{pacheco} Pacheco, Peter. \textit{An Introduction to Parallel Programming.} 2nd Edition, Morgan Kaufmann, 2021.
\bibitem{intel} Intel Corporation. \textit{Intel 64 and IA-32 Architectures Optimization Reference Manual.} Order Number: 248966-046A, 2023.
\bibitem{postgres} PostgreSQL Global Development Group. \textit{PostgreSQL 18.0 Documentation.} 2024.
\bibitem{drepper} Drepper, Ulrich. "What Every Programmer Should Know About Memory." \textit{Red Hat, Inc.}, 2007.
\bibitem{chandra} Chandra, Rohit, et al. \textit{Parallel Programming in OpenMP.} Morgan Kaufmann Publishers, 2001.
\end{thebibliography}

\end{document}
